% Clase 26 --- Dispersion cromatica en un prisma (§1.2).
% La luz blanca se abre porque n depende de lambda: cada color obedece la ley de
% Snell con SU propio indice y sale con otro angulo.
\begin{center}
\begin{tikzpicture}[scale=1]

  % prisma
  \fill[figblue!10] (0,0) -- (3,0) -- (1.5,2.6) -- cycle;
  \draw[figline, line width=1pt] (0,0) -- (3,0) -- (1.5,2.6) -- cycle;
  \node[figlblaux, anchor=north] at (1.5,-0.1) {prisma de vidrio};

  % haz blanco incidente
  \draw[figgray, line width=1.6pt, -{Latex[length=2.4mm,width=1.8mm]}]
    (-2.3,1.72) -- (0.55,1.0);
  \node[figlbl, anchor=south] at (-1.5,1.62) {luz blanca};

  % rayos internos (ya separados, apenas)
  \draw[violet!75, line width=1pt] (0.6,1.04) -- (2.13,1.51);
  \draw[green!55!black, line width=1pt] (0.6,1.04) -- (2.19,1.40);
  \draw[red!80, line width=1pt] (0.6,1.04) -- (2.25,1.29);

  % rayos emergentes, abiertos en abanico
  \draw[violet!75, line width=1pt, -{Latex[length=2.2mm,width=1.6mm]}]
    (2.13,1.51) -- (5.4,0.42);
  \draw[blue!60, line width=1pt, -{Latex[length=2.2mm,width=1.6mm]}]
    (2.15,1.47) -- (5.4,0.66);
  \draw[green!55!black, line width=1pt, -{Latex[length=2.2mm,width=1.6mm]}]
    (2.19,1.40) -- (5.4,0.92);
  \draw[yellow!75!black, line width=1pt, -{Latex[length=2.2mm,width=1.6mm]}]
    (2.22,1.35) -- (5.4,1.16);
  \draw[orange!85, line width=1pt, -{Latex[length=2.2mm,width=1.6mm]}]
    (2.24,1.31) -- (5.4,1.38);
  \draw[red!80, line width=1pt, -{Latex[length=2.2mm,width=1.6mm]}]
    (2.25,1.29) -- (5.4,1.60);

  \node[figlbl, violet!75, anchor=west] at (5.5,0.42) {violeta};
  \node[figlbl, red!80, anchor=west] at (5.5,1.60) {rojo};
  \node[figlblaux, anchor=north west] at (2.9,-0.35)
    {$n_{\text{violeta}} > n_{\text{rojo}}$: el violeta se desvía más};

\end{tikzpicture}\\[0.5em]
{\small\itshape La luz blanca es la que trae todos los colores en proporciones
parecidas, y el prisma no los \emph{crea}: los \textbf{separa}. La razón es que
el índice de refracción no es un número del material sino una función de la
longitud de onda, de modo que cada color entra con el mismo $\theta_1$ pero sale
con su propio $\theta_3$. El violeta, con $n$ mayor, es el que más se desvía. El
efecto se acumula en las dos caras del prisma, y por eso el abanico se ve
claramente aunque los índices difieran apenas en centésimas. Es exactamente el
mismo mecanismo del \textbf{arcoíris} ---la luz del Sol refractándose en las
gotitas de agua en suspensión--- y el que \enquote{mancha} de colores el borde
de un objeto visto a través de un vaso con agua. La moraleja operativa es que la
ley de Snell se aplica \textbf{longitud de onda por longitud de onda}: hablar
del ángulo refractado de la luz blanca, sin más, no tiene sentido.}
\end{center}
